Der Daten-Fit beschreibt die Passung zwischen den Datenanforderungen eines ERP-Systems 
und der vorhandenen Datenbasis eines KMU. Da ERP-Systeme auf einer zentralen oder gemeinsamen 
Datenbasis beruhen, müssen relevante Unternehmensdaten so vorliegen, dass sie in das System überführt und 
dort bereichsübergreifend genutzt werden können. Für die Bewertung 
des Daten-Fits ist daher entscheidend, ob zentrale Daten ausreichend vollständig, korrekt, aktuell und einheitlich sind. 

Dies betrifft insbesondere Bestands- und Inventurdaten, Verkaufs- und Kundendaten, Stücklisten, Personaldaten, 
Finanzdaten und Produktionskapazitäten \parencite{Jha2008,Xia2009,Alaskari2021,Leyh2017}. 
Sind diese Daten ungenau oder nicht systematisch gepflegt, entstehen Risiken für die Datenmigration, die Systemnutzung und die Qualität 
späterer Auswertungen \parencite{Jha2008,Leyh2017}. 

Die systematische Datenerfassung in KMU ist oft gering ausgeprägt. Daten, die für den Betrieb
eines ERP-Systems benötigt werden, fehlen häufig. Viele Daten sind
möglicherweise gar nicht vorhanden oder nicht korrekt \parencite{Jha2008}. Daher ist es wichtig,
Methoden zur Datenerhebung aufzubauen und Verantwortlichkeiten festzulegen, um
geeignete beziehungsweise verwertbare Daten zu erhalten.

Auf dieser Grundlage lässt sich Daten-Fit nicht nur als allgemeine
Datenqualität verstehen, sondern als konkrete Einsatzbereitschaft der Daten
für eine ERP-Einführung. Daten sind dann ERP-bereit, wenn sie ausreichend
genau, konsistent, bereinigt und strukturell in das Zielsystem überführbar
sind \parencite{Jha2008,Leyh2017}. Besonders bei zentralen Stammdaten, Bestandsdaten und Stücklisten ist
eine hohe Genauigkeit erforderlich, da fehlerhafte Daten die Nutzbarkeit des
ERP-Systems für Planung, Steuerung und Kontrolle erheblich beeinträchtigen
können \parencite{Jha2008}.

Ein wichtiges Kriterium ist die Fehlerfreiheit der Daten. Für Bestandsdaten und
Stücklisten wird eine Genauigkeit von mindestens 98\,\% genannt, damit das
ERP-System sinnvoll zur Unternehmenssteuerung eingesetzt werden kann
\parencite{Jha2008}. Liegt die Datenqualität deutlich darunter, besteht die
Gefahr, dass Prozesse zwar technisch im System abgebildet werden, die darauf
basierenden Informationen aber keine verlässliche Entscheidungsgrundlage
darstellen \parencite{Jha2008,Leyh2015}.

Neben der Genauigkeit ist auch die Bereinigung der vorhandenen Datenbestände
entscheidend. Fehlerhafte, veraltete oder korrupte Datensätze müssen vor der
Migration identifiziert und entfernt werden \parencite{Alaskari2021}. Ebenso sind doppelte Datensätze
zu bereinigen, die insbesondere durch Insellösungen, Excel-Dateien oder
abteilungsspezifische Systeme entstehen können \parencite{Alaskari2021}. Erst wenn solche Redundanzen
reduziert und die Datenbestände konsolidiert wurden, können sie sinnvoll in
eine gemeinsame ERP-Datenbasis überführt werden.

Darüber hinaus müssen Datenformate und Klassifikationen vereinheitlicht
werden \parencite{Leyh2015}. Wenn verschiedene Abteilungen unterschiedliche Bezeichnungen,
Strukturen oder Klassifizierungsmuster verwenden, entstehen
Integrationsprobleme bei der Überführung in das ERP-System \parencite{Alaskari2021}. Daten gelten daher
erst dann als einsatzbereit, wenn diese unterschiedlichen Strukturen
harmonisiert und in ein konsistentes Format überführt wurden.

Die technische Einsatzbereitschaft der Daten kann zusätzlich durch Testläufe
geprüft werden. Dabei wird eine ausgewählte Stichprobe, etwa ein sogenanntes
Gold Sample, in das neue System übertragen, um Mapping-Regeln und
Feldzuweisungen zu kontrollieren. Daten sind erst dann als migrationsbereit
anzusehen, wenn die dabei auftretenden Fehler behoben sind und sichergestellt
ist, dass relevante Informationen den richtigen Softwarefeldern zugeordnet
werden können \parencite{Alaskari2021}.

Schließlich ist Daten-Fit nicht nur eine technische, sondern auch eine
organisatorische Voraussetzung. Gerade in KMU müssen Methoden zur
kontinuierlichen Datenerhebung und Datenpflege etabliert werden. Zudem müssen
Verantwortlichkeiten für Datenqualität klar zugewiesen sein \parencite{Xia2009,Jha2008}. Ohne solche
Zuständigkeiten besteht die Gefahr, dass auch zunächst bereinigte Daten nach
dem Go-Live erneut veralten oder uneinheitlich gepflegt werden.

Zusammenfassend liegt ein hoher Daten-Fit vor, wenn relevante Unternehmensdaten
bereinigt, einheitlich strukturiert, ausreichend genau und durch Testläufe
validiert sind. Für KMU bedeutet dies, dass die Datenbasis vor der
ERP-Einführung nicht vollständig perfekt sein muss, aber ein Mindestmaß an
Genauigkeit, Konsistenz und organisatorischer Pflege aufweisen muss. Nur dann
kann das ERP-System seine Funktion als zentrale Datenbasis und
entscheidungsunterstützendes System wirksam erfüllen.


